\documentclass{report}
\usepackage{amsmath,amssymb}
\setlength{\parindent}{0mm}
\setlength{\parskip}{1em}
\begin{document}
\begin{center}
\rule{6in}{1pt} \
{\large Kevin Carlberg \\
{\bf Applying model reduction to Krylov-subspace recycling: the POD-augmented conjugate-gradient method}}

7011 East Ave \\ MS 9159 \\ Livermore \\ CA 94550
\\
{\tt ktcarlb@sandia.gov}\\
Virginia Forstall\\
Ray Tuminaro\end{center}

This talk presents a new Krylov-subspace-recycling method for efficiently
solving sequences of linear systems of equations characterized by varying
right-hand sides and symmetric-positive-definite matrices. As opposed to
typical truncation strategies used in recycling such as deflation, we propose
a truncation method based on a technique from nonlinear model reduction:
goal-oriented proper orthogonal decomposition (POD).

This idea is inspired by the observation that model reduction aims to compute a
low-dimensional subspace that contains an \textit{accurate} solution; as such,
we expect the proposed method to generate a low-dimensional subspace that is
well suited for computing solutions that can satisfy \textit{inexact}
tolerances. In particular, we propose specific goal-oriented POD `ingredients'
that align the optimality properties of POD with the objective of
Krylov-subspace recycling.

To efficiently compute solutions in the resulting `POD-augmented' Krylov
subspace, we propose a novel hybrid direct/iterative three-stage method that
leverages 1) the optimal ordering of POD basis vectors, and 2)
well-conditioned reduced matrices. Numerical experiments performed on
solid-mechanics problems highlight the benefits of the proposed method over
existing approaches for Krylov-subspace recycling.


\end{document}
